\section{Roman Numeral to Integer}
\label{sec:str-roman-to-integer}

\problemstatement{Convert a Roman numeral string (using symbols
I, V, X, L, C, D, M) to its integer value.}

\begin{patternbox}
Roman numerals are almost additive, with a \textbf{subtractive exception}:
a smaller symbol before a larger one is subtracted (IV $=4$). So scan
left to right and, at each symbol, \emph{subtract} its value if the next
symbol is larger, otherwise add it.
\end{patternbox}

\subsection*{Add, unless the next is larger}

Map each symbol to its value. Walk the string; for each symbol, compare its
value to the next symbol's value. If the current value is smaller than the
next, it is part of a subtractive pair (like IX), so subtract it; otherwise
add it. The last symbol is always added. This one pass handles both additive
and subtractive forms without special-casing each pair.

\begin{ideabox}[A Smaller Symbol Before a Larger One Subtracts]
The only irregularity in Roman numerals is the subtractive pair. Detecting
it locally -- ``current value $<$ next value'' -- converts the whole numeral
in a single left-to-right pass.
\end{ideabox}

\subsection*{Algorithm}

\begin{enumerate}
  \item Map symbols to values.
  \item For each position $i$: if \texttt{val[i] < val[i+1]}, subtract
        \texttt{val[i]}, else add it.
  \item Return the accumulated total.
\end{enumerate}

\subsection*{Dry run}

\dryrun{\texttt{"MCMXCIV"} ($1994$).}

\begin{itemize}
  \item M$+1000$; C$<$M so $-100$; M$+1000$; X$<$C so $-10$; C$+100$;
        I$<$V so $-1$; V$+5$.
  \item Total $=1000-100+1000-10+100-1+5=1994$.
\end{itemize}

\subsection*{C++ implementation}

\begin{lstlisting}
#include <bits/stdc++.h>
using namespace std;

int romanToInt(const string& s) {
    unordered_map<char,int> val = {
        {'I',1},{'V',5},{'X',10},{'L',50},{'C',100},{'D',500},{'M',1000}
    };
    int total = 0, n = s.size();
    for (int i = 0; i < n; i++) {
        if (i + 1 < n && val[s[i]] < val[s[i + 1]]) total -= val[s[i]];
        else total += val[s[i]];
    }
    return total;
}

int main() {
    cout << romanToInt("MCMXCIV") << "\n";         // 1994
    cout << romanToInt("LVIII") << "\n";           // 58
    return 0;
}
\end{lstlisting}

\begin{complexitybox}
\textbf{Time Complexity.} $O(n)$. \textbf{Space Complexity.} $O(1)$ (a
fixed symbol table).
\end{complexitybox}

\begin{takeawaybox}
Roman-to-integer is additive with one subtractive rule, detected by
comparing each symbol to its successor. Handling the exception locally
avoids enumerating every subtractive pair -- a clean parsing pattern.
\end{takeawaybox}
